problem:  
Let $\mathcal{O}$ be a compact 4‑dimensional Riemannian orbifold with positive sectional curvature. Assume that each singular point $p_i$ of $\mathcal{O}$ is modeled on $\mathbb{C}^2 / \Gamma_i$, where $\Gamma_i \subset \mathrm{SU}(2)$ is a non‑cyclic finite subgroup of the binary polyhedral type (so that each local deformation corresponds to an ADE configuration of $(-2)$‑spheres). Let $\widetilde{M}$ denote a smooth manifold obtained by replacing small neighborhoods of each singularity by the corresponding crepant resolution $X_i$ and smoothly gluing along their boundaries $S^3 / \Gamma_i$.  

**Question / Conjecture:**  
Is it possible for $\widetilde{M}$ to admit a Riemannian metric of positive sectional curvature?  

Equivalently, let $Q_{\mathrm{ADE}}$ be the negative‑definite intersection form on $H_2(X_i;\mathbb{Z})$ induced by the exceptional divisor—combinatorially the root lattice of type $A\!D\!E$.  
Do there exist any simply connected, compact, smooth 4‑manifolds with positive sectional curvature whose intersection form contains a sublattice isomorphic to $Q_{\mathrm{ADE}}$?  

If not, can one prove that the presence of such a negative‑definite ADE configuration of embedded 2‑spheres constitutes a topological obstruction to positive sectional curvature—analogous to how certain intersection‑form restrictions (as in Freedman, Donaldson, and Hitchin‑Thorpe‑type inequalities) obstruct Einstein or scalar curvature metrics?  

Why is it a "good" problem:  
This formulation isolates a single, sharp conjectural barrier between **topology** and **geometry**. Whether the presence of an ADE configuration of $(-2)$‑spheres precludes positive sectional curvature cuts to the heart of how 4‑manifold topology constrains metric geometry.  

- **Simplicity:**  The question can be stated in one line—*can a smooth 4‑manifold with an ADE lattice in its intersection form ever have positive sectional curvature?*—yet its resolution would require deep new insight.  
- **Bridging Fields:**  It unites Riemannian geometry (positivity of sectional curvature), algebraic surface theory (crepant resolutions and ADE lattices), and 4‑manifold topology (intersection forms, Donaldson theory).  
- **Mathematical Depth:**  Positive sectional curvature is extremely rigid, and its known 4‑manifold examples (spheres, $\mathbb{C}P^2$, certain biquotients) have simple intersection forms. Introducing an ADE lattice challenges this rigidity and asks whether known topological constraints are already at their optimal boundary.  
- **Potential Impact:**  A proof of impossibility would reveal new topological obstructions for positive curvature, sharpening our understanding of the landscape of possible 4‑manifolds with strict curvature bounds. Conversely, a counterexample would revolutionize the known families of positively curved manifolds.  

The problem’s conceptual core is thus transparent, its mathematical content genuinely frontier‑level, and its outcome—either way—would alter the state of knowledge in global differential geometry.